\documentclass[11pt]{article}
\usepackage[a4paper,margin=25mm]{geometry}
\usepackage{amsmath,amssymb,mathtools}
\usepackage{fontspec,unicode-math}
\setmainfont{Latin Modern Roman}
\setsansfont{Latin Modern Sans}
\setmonofont{Latin Modern Mono}
\setmathfont{Latin Modern Math}
\usepackage{xurl,hyperref}
\hypersetup{colorlinks=true,urlcolor=blue,linkcolor=blue}
\newcommand{\sref}[2]{\hyperlink{source:#1:#2}{[S#2]}}
\newcommand{\cataloguescope}[1]{\paragraph{Mathematical question.}#1}
\begin{document}
\section{Fast construction of five-term geometric progressions for NFS}\label{statement}

\subsection{Definitions and mathematical statement}

\paragraph{Shared notation.}$\mathbb N=\{1,2,\ldots\}$, and $\mathbb Z,\mathbb Q,\mathbb R,\mathbb C$ denote the integers, rationals, reals and complex numbers. Any other conventions are specified locally.


A five-term geometric progression modulo an integer $N>1$ is an integer tuple $(r_0,\ldots,r_4)$ with a ratio $s\in\mathbb Z/N\mathbb Z$ such that $r_j\equiv r_0s^j\pmod N$ for $0\le j\le4$. A rational linear recurrence of order at most two means a relation $a r_{j+2}+b r_{j+1}+c r_j=0$ for $j=0,1,2$, for rational $a,b,c$ not all zero. The source's usable progression also has $r_0$ invertible modulo $N$, so the ratio $r_1/r_0$ is defined.
\paragraph{Statement recorded in the supplied source.}
Find a practical method, for large $N$, to construct such progressions with
\[
\max_{0\le j\le4}|r_j|=O(N^{2/3})
\]
which satisfy no rational linear recurrence of order at most two. Use the associated two-dimensional orthogonal lattice to obtain two coprime cubic polynomials in $\mathbb Z[X]$ with a common root modulo $N$ and coefficients of the targeted order $O(N^{1/6})$. The cubic-size prediction in the slides is a typical-size goal, rather than an established worst-case guarantee. \sref{statement}{2}

\subsection{Short English statement}
Efficiently find short five-term modular geometric progressions that avoid low-order rational recurrences and yield small coprime cubic polynomials for the number field sieve.
\subsection{Sources}
\begin{description}
\item[\hypertarget{source:statement:1}{S1}] ULAM.AI and the UnsolvedMath Contributors, UnsolvedMath: A Curated Collection of Open Mathematics Problems, catalogue entry. CC BY 4.0. \url{https://huggingface.co/datasets/ulamai/UnsolvedMath}.
\item[\hypertarget{source:statement:2}{S2}] Peter L. Montgomery, \emph{Searching for Higher-Degree Polynomials for the General Number Field Sieve}, IPAM Number Theory and Cryptography: Open Problems workshop (9--13 October 2006), slides headed ``Five-term GP → Two Cubics and ``Desire a five-term GP mod N. \url{https://helper.ipam.ucla.edu/publications/SCWS1/SCWS1_6223.ppt}.
\end{description}
\label{end:statement}
\end{document}
